% SecondChance -- Project Synopsis
% Compile on Overleaf (overleaf.com) -- New Project > Upload Project, or paste
% this file in. No local LaTeX installation needed.
\documentclass[11pt,a4paper]{article}

\usepackage[margin=2.4cm]{geometry}
\usepackage[T1]{fontenc}
\usepackage{lmodern}
\usepackage{microtype}
\usepackage{booktabs}
\usepackage{array}
\usepackage{enumitem}
\usepackage{titlesec}
\usepackage{xcolor}
\usepackage[hidelinks]{hyperref}
\usepackage{graphicx}

\definecolor{jiitblue}{RGB}{23,54,93}

\titleformat{\section}{\color{jiitblue}\normalfont\Large\bfseries}{}{0pt}{}[\vspace{-6pt}\rule{\textwidth}{0.6pt}]
\titleformat{\subsection}{\normalfont\large\bfseries}{\thesubsection.}{0.5em}{}
\setlist{nosep,leftmargin=1.4em}
\renewcommand{\arraystretch}{1.25}

\newcolumntype{L}[1]{>{\raggedright\arraybackslash}p{#1}}

\begin{document}
\pagestyle{empty}

%% ---------------------------------------------------------------- title page
\begin{center}
\vspace*{1.5cm}

% Uncomment once you add the college logo to the project folder:
% \includegraphics[width=0.28\textwidth]{jiit-logo.png}\\[1.2cm]

{\color{jiitblue}\Large\bfseries JAYPEE INSTITUTE OF INFORMATION TECHNOLOGY}\\[4pt]
{\large Noida, Uttar Pradesh}\\[2.2cm]

{\large P\,R\,O\,J\,E\,C\,T\quad S\,Y\,N\,O\,P\,S\,I\,S}\\[4pt]
{\large Machine Learning and Data Mining}\\[1.0cm]

\rule{0.62\textwidth}{0.8pt}\\[0.7cm]
{\color{jiitblue}\Huge\bfseries SecondChance}\\[0.7cm]
{\large Mining Attribute Bias in Animal Shelter Outcomes with\\
Association Rules and Density-Based Clustering}\\[0.6cm]
\rule{0.62\textwidth}{0.8pt}

\vfill

\begin{minipage}[t]{0.45\textwidth}
\raggedright
{\color{jiitblue}\bfseries Submitted to:}\\[6pt]
Dr. Ashish Mishra
\end{minipage}
\hfill
\begin{minipage}[t]{0.45\textwidth}
\raggedleft
{\color{jiitblue}\bfseries Submitted by:}\\[6pt]
\begin{tabular}{@{}lr@{}}
Sambhav Ghildiyal & 23803005 \\
Aniket Pandey & 23803008 \\
Aditya Bansal & 23803009 \\
Mudit Maheshwari & 23803017 \\
\end{tabular}
\end{minipage}

\vspace{1.4cm}
November 2026
\end{center}

\newpage
\pagestyle{plain}
\setcounter{page}{1}

%% ------------------------------------------------------------------ abstract
\section*{Abstract}
\rule{\textwidth}{0.6pt}\vspace{4pt}

Municipal animal shelters decide, under permanent capacity pressure, which animals
are promoted for adoption, transferred to partner organisations, returned to owners,
or euthanised. Practitioners hold strong beliefs about what determines an animal's
fate: that black animals are passed over, that bully-breed labels are disqualifying,
that senior animals are effectively unadoptable. These beliefs already shape policy,
from shelter photography guidelines to the removal of breed labels from kennel cards.
They are rarely quantified at scale.

\textbf{SecondChance} is an end-to-end data mining study of \textbf{173{,}775 outcome
records} published by Austin Animal Center. The project will (i) mine association
rules over animal attributes to surface interpretable outcome patterns; (ii) benchmark
five classification algorithms predicting an animal's outcome under a leakage-safe
temporal split; (iii) segment the population with five clustering algorithms, using
dissimilarity measures appropriate to categorical data rather than defaulting to
Euclidean distance; (iv) test whether animals that density-based clustering labels as
noise --- that is, atypical animals --- experience measurably worse outcomes; and
(v) quantify the marginal effect of colour, breed group and age on adoption while
controlling for the remaining attributes. The deliverables are a reproducible pipeline,
a comparative study of the classical data mining toolkit on a large categorical
dataset, and an interpretable outcome-risk tool that returns a prediction together
with the rule that produced it.

%% --------------------------------------------------------------- motivation
\section*{Problem Statement and Motivation}

Shelter intake exceeds adoption capacity almost everywhere, so triage is unavoidable
and its consequences are terminal. Staff decide which animals receive foster
placement, promotion or extended holding time, and those decisions are guided largely
by experience and folklore rather than by measurement. The widely cited ``black dog
syndrome'' --- the claim that dark-coated animals wait longer and are adopted less
often --- is a representative example: it is believed strongly enough to drive
operational policy, yet reported findings in the literature are mixed, and most studies
are based on a single shelter over a short window.

A public dataset of 173{,}775 outcomes over more than a decade makes the question
tractable. Two questions are open:

\begin{enumerate}[label=\textbf{\arabic*.}]
  \item \textbf{Which attributes actually determine an animal's outcome, and by how
        much}, once the remaining attributes are controlled for?
  \item \textbf{Can an adverse outcome be anticipated early enough to act on it}, using
        only the information recorded when the animal enters the shelter?
\end{enumerate}

Interpretability is a requirement rather than a preference here. A shelter manager
cannot act on a risk score; they need the reason behind it in order to change
something --- a photograph, a kennel card, a transfer request. This makes decision
trees and association rules the appropriate instruments for the task, and makes a
comparison against less transparent models a genuine trade-off to measure rather than
a formality.

%% ---------------------------------------------------------------- objectives
\section*{Objectives}

\begin{enumerate}[label=\arabic*.]
  \item Build a reproducible pipeline and exploratory analysis: outcome distribution,
        trends over time, and normalisation of 3{,}010 free-text breed strings and
        663 colour strings into analysable categories.
  \item Mine association rules over animal attributes using Apriori, ranking patterns
        by lift rather than confidence alone, and report the rule sets that
        characterise each outcome class.
  \item Benchmark Decision Tree, Naive Bayes, $k$-Nearest Neighbours, Support Vector
        Machine and a Neural Network under one protocol, reporting macro-averaged F1
        because the outcome classes are strongly imbalanced.
  \item Segment the population with K-Means, K-Medoids, Hierarchical, DBSCAN and
        OPTICS, using a categorical dissimilarity measure, and compare internal and
        external cluster validity.
  \item Test the \emph{atypicality hypothesis}: are animals labelled as noise by
        DBSCAN and OPTICS over-represented among adverse outcomes?
  \item Quantify the marginal effect of colour, breed group, age and sterilisation
        status on adoption probability, and deliver an interpretable outcome-risk tool.
\end{enumerate}

%% ------------------------------------------------------------------- dataset
\section*{Dataset}

\begin{center}
\begin{tabular}{@{}L{3.1cm}L{12.2cm}@{}}
\toprule
\textbf{Property} & \textbf{Austin Animal Center} \\
\midrule
Outcomes & 173{,}775 outcome events, Oct 2013 -- May 2025, updated continuously \\
Intakes & 173{,}812 intake events; joins to Outcomes on Animal ID at
\textbf{99.1\,\%} coverage \\
Analysed & 162{,}461 dog and cat outcome events after scoping \\
Features & animal type, breed (3{,}010 free-text strings), colour (663), sex and
sterilisation status, age, intake type, intake condition \\
Label & Outcome --- Adoption 83{,}866; Transfer 46{,}668; Return to Owner
25{,}466; Euthanasia 4{,}704; Other 1{,}757 \\
Access & Two CSVs from the City of Austin open data portal, 21\,MB and 26\,MB,
no account or API key required \\
\bottomrule
\end{tabular}
\end{center}

\vspace{4pt}
The data is small enough to process entirely on a laptop CPU; no GPU is required at
any stage. The complete pipeline reproduces every table and figure in the report in
under ten minutes.

\subsection*{Three properties of the data that shape the method}

\textbf{Mixed timezone offsets.} The \texttt{DateTime} column carries mixed offsets
in both files, so it must be parsed as UTC; otherwise naive and aware timestamps
cannot be compared at all.

\textbf{Animals return.} About 10\,\% of Animal IDs appear more than once in the
Intakes file. The join must therefore be per \emph{outcome event}, matched to the
most recent intake at or before it, rather than one row per animal.

\textbf{The obvious feature leaks the label.} \texttt{Sex upon Outcome} encodes
sterilisation status, and shelters routinely sterilise animals \emph{as part of}
preparing them for adoption. Using it as a predictor is circular. All features are
therefore taken from the \textbf{intake} record, which is also what predicting at
intake honestly means.

%% --------------------------------------------------------------- methodology
\section*{Proposed Methodology}

\begin{center}
\small
\begin{tabular}{@{}cccccc@{}}
\toprule
\textbf{1. Ingest \& EDA} & $\rightarrow$ & \textbf{2. Rule Mining} & $\rightarrow$ &
\textbf{3. Classification} & \\
breed/colour normalisation & & Apriori, lift ranking & & five models, temporal split & \\
\midrule
\textbf{4. Clustering} & $\rightarrow$ & \textbf{5. Bias \& Tool} & & & \\
categorical dissimilarity & & marginal effects, risk tool & & & \\
\bottomrule
\end{tabular}
\end{center}

\subsection*{1. Preprocessing and Feature Engineering}
Breed strings are normalised into approximately twenty groups with a separate
mixed-breed indicator; colour strings are decomposed into a primary colour and a
pattern; age at outcome is parsed from free text into days; the combined sex field is
split into sex and sterilisation status, which are distinct predictors. Outcome month
and season are derived, since shelter intake is strongly seasonal. Every normalisation
rule is recorded in a mapping table reproduced in the report, so the transformation is
auditable rather than implicit.

\subsection*{2. Association Rule Mining}
Each animal is represented as a transaction over its attributes, for example
$\{$\texttt{dog}, \texttt{pit\_bull\_mix}, \texttt{black}, \texttt{adult},
\texttt{intact\_male}$\}$. Apriori is applied at a range of minimum support
thresholds, selected from the item frequency distribution rather than chosen
arbitrarily. Rules are ranked by lift, since confidence alone is misleading when one
outcome dominates. The resulting rule sets are compared directly against the paths
learned by the decision tree, which is an independent route to the same kind of
structure.

\subsection*{3. Classification}
Decision Tree, Gaussian and Categorical Naive Bayes, $k$-Nearest Neighbours, Support
Vector Machine (linear and RBF) and a feedforward neural network trained with
backpropagation are evaluated on an identical test set. Categorical attributes are
one-hot encoded; class weighting addresses imbalance. Evaluation uses stratified
$k$-fold cross-validation together with a \textbf{temporal holdout} --- training on
earlier years and testing on later ones --- because a random split would allow the
model to learn period-specific shelter policy and report an optimistic result.

\subsection*{4. Clustering with Categorical Dissimilarity}
The attributes are categorical, so Euclidean distance is not meaningful and the choice
of dissimilarity is a methodological decision rather than a detail. K-Medoids,
Hierarchical, DBSCAN and OPTICS all accept an arbitrary precomputed dissimilarity
matrix and are run with Hamming and Jaccard measures. K-Means cannot: it computes
means, which requires Euclidean geometry, and the mean of categorical vectors is not
an animal. K-Means on one-hot encoded data is therefore included deliberately as the
contrast that demonstrates why medoid-based clustering exists. K-Medoids is
implemented directly, since \texttt{scikit-learn-extra} is unmaintained and does not
build on current Python versions.

\subsection*{5. Bias Quantification and Outcome-Risk Tool}
The marginal effect of each attribute on adoption probability is estimated while
holding the remaining attributes fixed, which separates genuine attribute effects from
confounding --- black animals are disproportionately common among certain breed groups,
so an unadjusted comparison measures the wrong thing. The final deliverable is a tool
that accepts an animal's attributes and returns an outcome risk together with the
association rule or decision path responsible for it.

%% ------------------------------------------------------------------- metrics
\section*{Evaluation Metrics}

\textbf{Macro-averaged F1} is the primary classification metric, because accuracy is
uninformative when one outcome accounts for roughly half of all records. Per-class
precision and recall, balanced accuracy and the Matthews Correlation Coefficient are
also reported. Association rules are assessed by support, confidence and \textbf{lift}.
Clustering is assessed externally by Adjusted Rand Index and Normalised Mutual
Information against the outcome labels, internally by silhouette computed on the
dissimilarity matrix, and by the number of clusters that are near-pure in a single
outcome class. All results are reported as mean $\pm$ standard deviation over five
random seeds, alongside the temporal holdout result.

%% --------------------------------------------------------------------- tools
\section*{Tools and Technology Stack}

\begin{center}
\begin{tabular}{@{}L{3.6cm}L{11.7cm}@{}}
\toprule
\textbf{Layer} & \textbf{Tools} \\
\midrule
Language \& environment & Python 3.13, Jupyter, Git / GitHub \\
Data processing & pandas, NumPy \\
Association rules & mlxtend (Apriori, FP-Growth) \\
Machine learning & scikit-learn (trees, Naive Bayes, KNN, SVM, MLP, clustering) \\
Neural network & PyTorch --- categorical embeddings and a class-weighted loss,
neither of which \texttt{MLPClassifier} supports; CPU only \\
Clustering support & SciPy (linkage, dendrograms), custom K-Medoids implementation \\
Visualisation & Matplotlib \\
Application & Streamlit outcome-risk tool \\
Hardware & Laptop CPU throughout; no GPU required \\
\bottomrule
\end{tabular}
\end{center}

%% ----------------------------------------------------------------- work plan
\section*{Work Plan}

\begin{center}
\begin{tabular}{@{}cL{11.0cm}c@{}}
\toprule
\textbf{Phase} & \textbf{Activities} & \textbf{Duration} \\
\midrule
1 & Data ingestion, exploratory analysis, breed and colour normalisation, baseline
decision tree & Week 1 \\
2 & Association rule mining, support threshold selection, lift-ranked rule
interpretation & Week 2 \\
3 & Classification benchmark, imbalance handling, temporal holdout evaluation & Week 3 \\
4 & Clustering suite, categorical dissimilarity, atypicality hypothesis test & Week 4 \\
5 & Bias quantification, outcome-risk tool, report and presentation & Week 5 \\
\bottomrule
\end{tabular}
\end{center}

\vspace{6pt}
\noindent\textbf{Responsibilities.}
\emph{Sambhav Ghildiyal:} data pipeline, breed and colour normalisation, exploratory analysis
and all descriptive figures.
\emph{Aniket Pandey:} association rule mining, support and lift analysis, and the
comparison between mined rules and decision tree paths.
\emph{Aditya Bansal:} the five classification algorithms, imbalance handling and the
temporal evaluation protocol.
\emph{Mudit Maheshwari:} the five clustering algorithms, categorical dissimilarity measures
and the atypicality hypothesis.
All members share the bias quantification, the final report and the presentation.

%% ---------------------------------------------------------------- outcomes
\section*{Preliminary Results}

The pipeline is implemented and has been run end to end. Early findings:

\textbf{The folk belief does not survive measurement.} Black dogs are adopted
\emph{more} often than other dogs (+2.77 percentage points). Black cats are adopted
slightly less (--2.25 pp) and wait about one day longer. After direct
standardisation for species, age band and intake route, black sits mid-table and
\emph{white} animals fare worst. ``Black dog syndrome'' is not supported by these
173{,}775 records.

\textbf{What does matter is age, and then breed label.} Senior dogs are adopted
23.66 pp less often than non-seniors; senior cats 10.32 pp less. Pit bull type dogs
are adopted 4.34 pp less --- but wait a median of \textbf{10.0 days against 5.3}.
The trait that attracts public attention is close to noise; the traits that do not
are large.

\textbf{Atypical animals fare worse.} Animals that DBSCAN and OPTICS label as noise
--- unusual attribute combinations, identified with no labels whatsoever --- are
euthanised at roughly \textbf{2.7 to 3.2 times} the rate of animals that fall inside
a cluster, and are adopted less often. Labels are used only to score the finding,
never to produce it.

\textbf{Honest evaluation costs accuracy.} Every classifier degrades under the
temporal holdout relative to a random split. A random split lets a model learn
period-specific shelter policy and take credit for it.

\section*{Expected Outcomes}

\begin{itemize}
  \item A reproducible, open-source pipeline running from the raw public CSVs to a
        trained, interpretable outcome predictor.
  \item A quantified answer to a question shelters act on but rarely measure: the
        effect of colour, breed group and age on adoption, with confounders
        controlled by direct standardisation.
  \item A comparison of mined association rules against learned decision tree paths,
        establishing where two independent methods agree and, more informatively,
        where and why they do not.
  \item An empirical test of whether atypical animals, identified without labels by
        density-based clustering, experience worse outcomes.
  \item An outcome-risk tool that returns both a prediction and the rule responsible
        for it, usable in a shelter triage workflow.
\end{itemize}

%% -------------------------------------------------------------- references
\section*{References}

\begin{enumerate}[label={[\arabic*]}]
  \item City of Austin, ``Austin Animal Center Outcomes,'' City of Austin Open Data
        Portal. [Online]. Available: \url{https://data.austintexas.gov/}
  \item R. Agrawal and R. Srikant, ``Fast Algorithms for Mining Association Rules,''
        in \emph{Proc. 20th Int. Conf. Very Large Data Bases (VLDB)}, 1994,
        pp.~487--499.
  \item M. Ester, H.-P. Kriegel, J. Sander, and X. Xu, ``A Density-Based Algorithm for
        Discovering Clusters in Large Spatial Databases with Noise,'' in
        \emph{Proc. 2nd Int. Conf. Knowledge Discovery and Data Mining (KDD)}, 1996,
        pp.~226--231.
  \item M. Ankerst, M. M. Breunig, H.-P. Kriegel, and J. Sander, ``OPTICS: Ordering
        Points To Identify the Clustering Structure,'' in \emph{Proc. ACM SIGMOD Int.
        Conf. Management of Data}, 1999, pp.~49--60.
  \item L. Kaufman and P. J. Rousseeuw, \emph{Finding Groups in Data: An Introduction
        to Cluster Analysis}. New York: Wiley, 1990.
  \item Z. Huang, ``Extensions to the $k$-Means Algorithm for Clustering Large Data
        Sets with Categorical Values,'' \emph{Data Mining and Knowledge Discovery},
        vol.~2, no.~3, pp.~283--304, 1998.
  \item F. Pedregosa \emph{et al.}, ``Scikit-learn: Machine Learning in Python,''
        \emph{Journal of Machine Learning Research}, vol.~12, pp.~2825--2830, 2011.
  \item S. Raschka, ``MLxtend: Providing machine learning and data science utilities
        and extensions to Python's scientific computing stack,'' \emph{Journal of Open
        Source Software}, vol.~3, no.~24, p.~638, 2018.
\end{enumerate}

\vspace{6pt}
\noindent\footnotesize\emph{Note to the group: references [2]--[8] are the standard
primary sources for the algorithms and tools used. Before submission, add two or three
domain references on shelter adoption outcomes and ``black dog syndrome'' from the
animal-welfare literature (\emph{Anthrozo\"os} and the \emph{Journal of Applied Animal
Welfare Science} both publish on this) and verify every citation yourself.}

\end{document}
